% ============================================================
% methods.tex
% ============================================================

\section{Methods}
\label{sec:methods}

\subsection{Participants and Datasets}
\label{sec:methods-data}

Two publicly available datasets were used, chosen to span complementary segments of the adult lifespan and to provide well-validated behavioral phenotyping alongside high-resolution structural MRI.

\subsubsection{Human Connectome Project -- Young Adult (HCP-YA)}

The HCP-YA dataset~\cite{vanessen2013wuminn} comprises healthy adults aged 22 to 37 who were recruited as part of the WU-Minn consortium.
After excluding participants with missing values on any analysis variable (sleep quality, fluid intelligence, neuroticism, or FreeSurfer-derived structural metrics), the final analytic sample consisted of \textit{N}\,=\,1{,}104 individuals.
Participants completed the Pittsburgh Sleep Quality Index (PSQI;~\cite{buysse1989pittsburgh}), providing a global sleep-quality score with higher values indicating poorer sleep; the Penn Progressive Matrices (PMAT24\_A\_CR), a measure of fluid reasoning; and the NEO Five-Factor Inventory neuroticism subscale (NEOFAC\_N;~\cite{costa1992neo}).
Each participant's T1-weighted MRI was processed through the FreeSurfer pipeline~\cite{fischl2012freesurfer}, yielding estimates of regional cortical thickness and cortical surface area for 34 regions per hemisphere (68 total) based on the Desikan-Killiany (DK) atlas~\cite{desikan2006automated}.
Only baseline-visit scans were retained when multiple visits were available.
Age was recorded in bins (e.g., 22--25, 26--30, 31--35); we used the midpoint of each bin for descriptive purposes.

\subsubsection{Aging Brain and Behaviour Cohort (AABC Release 2)}

The AABC study~\cite{bookheimer2019aging} recruited a community-dwelling sample spanning ages 36 to 89, with data acquired across multiple sites as part of the Lifespan Human Connectome Project in Aging.
Structural MRI data were processed using FreeSurfer, and cortical parcellation followed the Glasser HCP-MMP1 atlas~\cite{glasser2016multimodal}, which subdivides each hemisphere into 180 cytoarchitectonically defined regions (360 total).
Behavioral data were obtained from the extended non-imaging battery, which includes the PSQI global sleep-quality score and the NEO neuroticism subscale.
After merging demographic, behavioral, and structural data and retaining only baseline visits, the final sample comprised \textit{N}\,=\,1{,}150 participants.
Age was recorded as a continuous variable, and for ANOVA analyses, participants were binned into decades (30s, 40s, 50s, 60s, 70s, 80+).

\subsection{Cortical Parcellation Atlases}
\label{sec:methods-atlases}

The two datasets employ different parcellation schemes, which differ in spatial resolution and in the principles underlying region definition.
The Desikan-Killiany atlas~\cite{desikan2006automated} defines 34 gyral-based regions per hemisphere on the basis of macroanatomical landmarks; it is widely used in clinical and population-level studies owing to its high test-retest reliability across scanners.
The Glasser HCP-MMP1 atlas~\cite{glasser2016multimodal} defines 180 regions per hemisphere using multimodal boundaries derived from cortical architecture, function, connectivity, and topography; it offers substantially finer granularity but requires high-quality data for accurate parcellation.
Because the two atlases partition the cortex at different spatial scales, region-level results cannot be compared directly.
To facilitate cross-dataset interpretation, we constructed a lobe-level correspondence table that maps each DK region and each Glasser region to one of five canonical lobes (Frontal, Parietal, Temporal, Occipital, Insular).
For DK regions in the cingulate cortex, anterior cingulate regions were assigned to the Frontal lobe and posterior/isthmus cingulate regions were assigned to the Parietal lobe, consistent with cytoarchitectonic conventions used in the Glasser atlas.
The resulting lobe-level mapping is visualized in Figure~\ref{fig:atlas-lobe}, which renders both atlases on the fsaverage inflated surface with shared lobe coloring.

\begin{figure}[htbp]
  \centering
  \includegraphics[width=0.95\textwidth]{atlas_lobe_correspondence.pdf}
  \caption{Lobe-level correspondence between the Desikan-Killiany atlas (left; 34 regions per hemisphere) and the Glasser HCP-MMP1 atlas (right; 180 regions per hemisphere), rendered on the fsaverage5 inflated surface.
           Both atlases are colored by the same five-lobe scheme (Frontal, Parietal, Temporal, Occipital, Insular).
           The central legend indicates the number of regions per lobe in each atlas.}
  \label{fig:atlas-lobe}
\end{figure}

\subsection{Statistical Analyses}
\label{sec:methods-stats}

All analyses were performed in Python 3.11 using \texttt{scipy.stats}, \texttt{statsmodels}, and \texttt{pandas}.
The analytic strategy was held constant across research questions to permit direct comparison of effect sizes.

\subsubsection{Group Comparisons}

For research questions involving a categorical grouping variable (sleep quality, neuroticism), participants were divided into two groups and each region of interest (ROI) was tested independently.
For sleep quality, participants scoring $\leq$\,5 on the PSQI were classified as good sleepers and those scoring $>$\,5 as poor sleepers~\cite{buysse1989pittsburgh}.
For neuroticism, the bottom and top tertiles of the NEOFAC\_N distribution were retained as the low- and high-neuroticism groups, respectively.
Between-group differences in cortical thickness (or surface area) at each ROI were assessed using Welch's \textit{t}-test, which does not assume equal variances.
The effect size for each comparison was quantified as Cohen's \textit{d} using the pooled standard deviation~\cite{cohen1988statistical}:
\begin{equation}
  d = \frac{\bar{X}_1 - \bar{X}_2}{\sqrt{\dfrac{(n_1 - 1)\,s_1^2 + (n_2 - 1)\,s_2^2}{n_1 + n_2 - 2}}}
  \label{eq:cohens-d}
\end{equation}
Ninety-five percent confidence intervals for \textit{d} were obtained using a normal approximation based on the standard error of \textit{d}~\cite{cohen1988statistical}.

\subsubsection{Correlational Analyses}

For research questions involving a continuous predictor (fluid intelligence, age), Pearson's product-moment correlation coefficient (\textit{r}) was computed between the predictor and the cortical metric at each ROI.
Ninety-five percent confidence intervals for \textit{r} were derived via the Fisher \textit{z}-transformation.
In the aging analyses, one-way ANOVA across age-decade bins was additionally computed for each FDR-significant region to characterize the functional form of the age effect and to test whether the decline was uniform across decades.

\subsubsection{Multiple-Comparison Correction}

Within each set of ROI-wise tests, \textit{p}-values were corrected for multiplicity using the Benjamini-Hochberg false discovery rate (FDR-BH) procedure~\cite{benjamini1995controlling} at \textit{q}\,<\,.05.
A region was considered statistically significant only if its FDR-corrected \textit{p}-value fell below .05.
This approach was chosen over family-wise error rate corrections (e.g., Bonferroni) because it provides a more favorable balance between Type~I and Type~II error rates in large-scale neuroimaging analyses where dozens to hundreds of tests are conducted simultaneously.

\subsubsection{Reporting Conventions}

All results are reported in APA format.
For \textit{t}-tests, we report \textit{t}(df)\,=\,\textit{value}, \textit{p}\,=\,\textit{value}, \textit{d}\,=\,\textit{value}, 95\% CI [\textit{lower}, \textit{upper}].
For correlations, we report \textit{r}(\textit{N})\,=\,\textit{value}, \textit{p}\,=\,\textit{value}, 95\% CI [\textit{lower}, \textit{upper}].
Exact \textit{p}-values are reported to three decimal places; values smaller than .001 are reported as \textit{p}\,<\,.001.
All effect sizes are accompanied by 95\% confidence intervals.
